\begin{question}

  This problem is a bit tricky.  To properly solve this problem,
  you'll need Real Analysis.  Specifically, you'll need to understand
  Distribution Theory.  Because the coordinate system itself breaks
  down at $r=0$ (where you no longer have the concept of orthogonality
  since $\varphi$ and $\theta$ don't really exist), the partial
  derivatives also break down.  Solving this sort of problem requires
  Distribution Theory to take Weak Derivatives (or something like that
  anyway).

  Putting a pin in this one since I don't have the requisite knowledge
  of Distribution Theory to solve this one at the moment.


  
The time-averaged potential of a neutral hydrogen atom is given by

\begin{equation}
  \Phi = \frac{q}{4\pi\epsilon_0} \frac{e^{-\alpha r}}{r} \left( 1 +
  \frac{\alpha r}{2} \right)
\end{equation}

where $q$ is the magnitude of the electronic charge, and $\alpha^-1 =
a_0/2$, $a_0$ being the Bohr radius. Find the distribution of charge
(both continuous and discrete) that will give this potential and
interpret your result physically.

\end{question}


We start with equation 1.28 from the book

\begin{equation}
  \nabla^2 \Phi = -\rho/\epsilon_0
\end{equation}

Since we're given $\Phi$, we can apply this relation to get $\rho$.

\begin{equation}
  \nabla^2 \Phi
  = -\rho/\epsilon_0
  = \nabla^2 \frac{q}{4\pi\epsilon_0} \frac{e^{-\alpha r}}{r} \left( 1 +
  \frac{\alpha r}{2} \right)
\end{equation}

Looking at the Laplacian in spherical coordinates, for the inside
cover, we find that the radial components are

\begin{equation}
  \nabla^2 \psi
  = \frac{1}{r^2}
  \frac{\partial}{\partial r}
  \left(r^2 \frac{\partial}{r}\psi \right)
\end{equation}

Applying that to our Electric potential we have

\begin{equation}
  -\frac{\rho}{\epsilon_0}
  = \frac{1}{r^2}
  \frac{\partial}{\partial r}
  \left\{
  r^2 \frac{\partial}{r}
  \frac{q}{4\pi\epsilon_0} \frac{e^{-\alpha r}}{r} \left( 1 +
  \frac{\alpha r}{2} \right)  
  \right\}
\end{equation}

Cleaning up a few things

\begin{equation}
  \rho
  = -\frac{q}{4\pi}
  \frac{1}{r^2}
  \frac{\partial}{\partial r}
  \left\{
  r^2 \frac{\partial}{r}
  \frac{e^{-\alpha r}}{r} \left( 1 + \frac{\alpha r}{2} \right)
  \right\}
\end{equation}

And diving right in

\begin{equation}
  \rho
  = -\frac{q}{4\pi}
  \frac{1}{r^2}
  \frac{\partial}{\partial r}
  \left\{
  r^2 \frac{\partial}{r}
  \frac{e^{-\alpha r}}{r}
  + r^2 \frac{\partial}{r}
  e^{-\alpha r}\frac{\alpha}{2}
  \right\}
\end{equation}

Evaluate the two derivatives

\begin{equation}
  \rho
  = -\frac{q}{4\pi}
  \frac{1}{r^2}
  \frac{\partial}{\partial r}
  \left\{
  r^2 \left(
  e^{-\alpha r}
  \frac{\partial}{\partial r}  
  \left(\frac{1}{r}\right)
  - \alpha \frac{e^{-\alpha r}}{r}
  \right)
  + r^2 \left(
  - \frac{\alpha^2}{2}
  e^{-\alpha r}
  \right)
  \right\}
\end{equation}

Separate the partial we can't take because spherical coordinates are
annoying.

\begin{equation}
  \rho
  = -\frac{q}{4\pi}
  \frac{1}{r^2}
  \frac{\partial}{\partial r}
  \left\{
  -\alpha r e^{-\alpha r}
  - \frac{r^2\alpha^2}{2}
  e^{-\alpha r}
  \right\}
  -\frac{q}{4\pi}  
  \frac{1}{r^2}
  \frac{\partial}{\partial r}
  \left\{
  r^2
  e^{-\alpha r}
  \frac{\partial}{\partial r}  
  \left(\frac{1}{r}\right)
  \right\}
\end{equation}

Cleaning up a bit

\begin{equation}
  \rho
  = \frac{q}{4\pi}
  \frac{1}{r^2}
  \frac{\partial}{\partial r}
  \left\{
  e^{-\alpha r}
  \left(
  \alpha r
  + \frac{r^2\alpha^2}{2}
  \right)
  \right\}
  -\frac{q}{4\pi}  
  \frac{1}{r^2}
  \frac{\partial}{\partial r}
  \left\{
  r^2
  e^{-\alpha r}
  \frac{\partial}{\partial r}  
  \left(\frac{1}{r}\right)
  \right\}
\end{equation}

Taking the next partial on the left

\begin{equation}
  \rho
  = \frac{q}{4\pi}
  \frac{1}{r^2}
  \left\{
  -\alpha e^{-\alpha r}
  \left(
  \alpha r
  + \frac{r^2\alpha^2}{2}
  \right)
  + e^{-\alpha r}
  \left(
  \alpha
  + \alpha^2 r
  \right)
  \right\}
  -\frac{q}{4\pi}  
  \frac{1}{r^2}
  \frac{\partial}{\partial r}
  \left\{
  r^2
  e^{-\alpha r}
  \frac{\partial}{\partial r}  
  \left(\frac{1}{r}\right)
  \right\}
\end{equation}

Distributing our terms

\begin{equation}
  \rho
  = \frac{q}{4\pi}
  \frac{1}{r^2}
  e^{-\alpha r}
  \left\{
  -\alpha^2 r
  - \frac{r^2\alpha^3}{2}
  + \alpha
  + \alpha^2 r
  \right\}
  -\frac{q}{4\pi}  
  \frac{1}{r^2}
  \frac{\partial}{\partial r}
  \left\{
  r^2
  e^{-\alpha r}
  \frac{\partial}{\partial r}  
  \left(\frac{1}{r}\right)
  \right\}
\end{equation}

Buttoning up the left side

\begin{equation}
  \rho
  = \frac{q}{4\pi}
  e^{-\alpha r}
  \left(
  - \frac{\alpha^3}{2}
  + \frac{\alpha}{r^2}
  \right)
  -\frac{q}{4\pi}  
  \frac{1}{r^2}
  \frac{\partial}{\partial r}
  \left\{
  r^2
  e^{-\alpha r}
  \frac{\partial}{\partial r}  
  \left(\frac{1}{r}\right)
  \right\}
\end{equation}

Because we've worked this problem already, we know that we need to
pull a Laplacian out of it so we can prevent it from becoming singular
at $r=0$.  To angle towards that, let's first take note of what the
Laplacian of $1/r$ looks like

\begin{equation}
  \nabla^2 \left(\frac{1}{r}\right)
  = \frac{1}{r^2} \frac{\partial}{\partial r}
  \left(r^2 \frac{\partial}{\partial r}\left(\frac{1}{r}\right) \right)
  = \frac{1}{r^2} \left(
  2r\frac{\partial}{\partial r}\left(\frac{1}{r}\right)
  + r^2\frac{\partial^2}{\partial r^2}\left(\frac{1}{r}\right)
  \right)
\end{equation}

So now let's try to make our equation look like that.  First, we'll
have a derivatve at hand for reference

\begin{equation}
  \frac{\partial}{\partial r}
  r^2
  e^{-\alpha r}
  = 2r\,e^{-\alpha r}
  - \alpha r^2\,e^{-\alpha r}
  = e^{-\alpha r} \left( 2r - \alpha r^2 \right)
\end{equation}

Using that formulation, we can take our last derivative to find

\begin{equation}
  \rho
  = \frac{q}{4\pi}
  e^{-\alpha r}
  \left(
  - \frac{\alpha^3}{2}
  + \frac{\alpha}{r^2}
  \right)
  -\frac{q}{4\pi}  
  \frac{1}{r^2}
  e^{-\alpha r}
  \left\{
  r^2
  \frac{\partial^2}{\partial r^2}(\frac{1}{r})
  + \left( 2r - \alpha r^2 \right)
  \frac{\partial}{\partial r}\left(\frac{1}{r}\right)
  \right\}
\end{equation}

Cleaning things up a bit, we find that

\begin{equation}
  \rho
  = \frac{q}{4\pi}
  e^{-\alpha r}
  \left(
  - \frac{\alpha^3}{2}
  + \frac{\alpha}{r^2}
  \right)
  -\frac{q}{4\pi}  
  e^{-\alpha r}
  \left[
  \frac{1}{r^2}
  \left\{
  r^2 \frac{\partial^2}{\partial r^2}(\frac{1}{r})
  + 2r \frac{\partial}{\partial r}\left(\frac{1}{r}\right)
  \right\}
  \right]
  -\frac{q}{4\pi}  
  e^{-\alpha r}
  \left(-\alpha \frac{\partial}{\partial r}\left(\frac{1}{r}\right) \right)
\end{equation}

This allows us to pull out the Laplacian in the middle...but we're not
going to do that...not yet.  First, let's evaluate that lone
derivative on the right.

\begin{equation}
  \rho
  = \frac{q}{4\pi}
  e^{-\alpha r}
  \left(
  - \frac{\alpha^3}{2}
  + \frac{\alpha}{r^2}
  \right)
  -\frac{q}{4\pi}  
  e^{-\alpha r}
  \left[
  \frac{1}{r^2}
  \left\{
  r^2 \frac{\partial^2}{\partial r^2}(\frac{1}{r})
  + 2r \frac{\partial}{\partial r}\left(\frac{1}{r}\right)
  \right\}
  \right]
  -\frac{q}{4\pi}  
  e^{-\alpha r}
  \frac{\alpha}{r^2}
\end{equation}

Noting that the newly evaluated term cancels out one on the left, we have

\begin{equation}
  \rho
  = - \frac{q}{8\pi}
  e^{-\alpha r}
  \alpha^3
  -\frac{q}{4\pi}  
  e^{-\alpha r}
  \left[
  \frac{1}{r^2}
  \left\{
  r^2 \frac{\partial^2}{\partial r^2}(\frac{1}{r})
  + 2r \frac{\partial}{\partial r}\left(\frac{1}{r}\right)
  \right\}
  \right]
\end{equation}

Now, instead of pulling out a laplacian, we're going to evaluate the
derivatives

\begin{equation}
  \rho
  = - \frac{q}{8\pi}
  e^{-\alpha r}
  \alpha^3
  -\frac{q}{4\pi}  
  e^{-\alpha r}
  \left[
  \frac{1}{r^2}
  \left\{
  r^2 \frac{\partial}{\partial r}(\frac{-1}{r^2})
  + 2r \frac{-1}{r^2}
  \right\}
  \right]
\end{equation}

Cleaning up and evaluating the final derivative, we get

\begin{equation}
  \rho
  = - \frac{q}{8\pi}
  e^{-\alpha r}
  \alpha^3
  -\frac{q}{4\pi}  
  e^{-\alpha r}
  \left[
  \frac{1}{r^2}
  \left\{
  r^2 (\frac{2}{r^3})
  - \frac{2}{r}
  \right\}
  \right]
\end{equation}

Which cleans up finally to

\begin{equation}
  \rho
  = - \frac{q}{8\pi}
  e^{-\alpha r}
  \alpha^3
\end{equation}








\begin{equation}
  \rho
  = \frac{q}{4\pi}
  e^{-\alpha r}
  \left(
  - \frac{\alpha^3}{2}
  + \frac{\alpha}{r^2}
  \right)
  -\frac{q}{4\pi}  
  e^{-\alpha r}
  \nabla^2 \left(\frac{1}{r}\right)
  -\frac{q}{4\pi}  
  e^{-\alpha r}
  \left(-\alpha \frac{\partial}{\partial r}\left(\frac{1}{r}\right) \right)
\end{equation}

Evaluating the final partial derivative and combining with the terms
on the left, we get

\begin{equation}
  \rho
  = \frac{q}{4\pi}
  e^{-\alpha r}
  \left(
  - \frac{\alpha^3}{2}
  + \frac{\alpha}{r^2}
  - \frac{\alpha}{r^2}
  \right)
  -\frac{q}{4\pi}  
  e^{-\alpha r}
  \nabla^2 \left(\frac{1}{r}\right)
\end{equation}

Since the partial derivative on the right cancels out the term on the
left, we can remove that and substitute in $-4\pi\delta(r)$ for our
Laplacian, leaving us with our final answer:

\begin{equation}
  \rho
  = q\delta(r)
  - \frac{q}{8\pi}
  e^{-\alpha r}
  \alpha^3
\end{equation}


Here we show a positive charge as a delta function at $r=0$
representing the proton in the Hydrogen atom followed by an
exponentially decaying spherical shell of negative charge for our
lonely little electron.
